\documentclass[10pt,a4paper]{article}

% Packages essentiels
\usepackage[utf8]{inputenc}
\usepackage[T1]{fontenc}
\usepackage[margin=1.5cm]{geometry}
\usepackage{amsmath,amssymb}
\usepackage{tikz}
\usetikzlibrary{positioning,calc,arrows.meta,shapes.geometric}
\usepackage{xcolor}
\usepackage{listings}
\usepackage{tcolorbox}
\tcbuselibrary{breakable}
\usepackage{hyperref}

% Configuration listings avec gestion complète des accents
\lstset{
    language=C,
    basicstyle=\ttfamily\scriptsize,
    backgroundcolor=\color{gray!10},
    keywordstyle=\color{blue},
    commentstyle=\color{gray},
    stringstyle=\color{red},
    numbers=left,
    numberstyle=\tiny,
    breaklines=true,
    frame=single,
    tabsize=4,
    showstringspaces=false,
    inputencoding=utf8,
    extendedchars=true,
    literate={é}{{\'e}}1 {è}{{\`e}}1 {à}{{\`a}}1 {ç}{{\c{c}}}1 
             {ê}{{\^e}}1 {î}{{\^i}}1 {ô}{{\^o}}1 {û}{{\^u}}1 {ù}{{\`u}}1
             {É}{{\'E}}1 {È}{{\`E}}1 {À}{{\`A}}1 {Ç}{{\c{C}}}1
             {Ê}{{\^E}}1 {Î}{{\^I}}1 {Ô}{{\^O}}1 {Û}{{\^U}}1 {Ù}{{\`U}}1
             {ë}{{\"e}}1 {ï}{{\"\i}}1 {ö}{{\"o}}1 {ü}{{\"u}}1
             {Ë}{{\"E}}1 {Ï}{{\"I}}1 {Ö}{{\"O}}1 {Ü}{{\"U}}1
             {œ}{{\oe}}1 {Œ}{{\OE}}1 {æ}{{\ae}}1 {Æ}{{\AE}}1
}

% Boxes pédagogiques
\newtcolorbox{exercicebox}[1]{
    colback=green!5,colframe=green!50!black,
    fonttitle=\bfseries,title=#1,breakable
}

\newtcolorbox{warningbox}[1]{
    colback=red!5,colframe=red!75!black,
    fonttitle=\bfseries,title=#1,breakable
}

\newtcolorbox{infobox}[1]{
    colback=blue!5,colframe=blue!75!black,
    fonttitle=\bfseries,title=#1,breakable
}

\newtcolorbox{techbox}[1]{
    colback=orange!5,colframe=orange!75!black,
    fonttitle=\bfseries,title=#1,breakable
}

\title{\textbf{TP Chapitre 3 : Structures et Unions}\\
\large Programmation C Avancée - L3}
\author{Dr. El Hadji Bassirou TOURÉ\\
Département de Mathématiques et Informatique\\
Faculté des Sciences et Techniques\\
Université Cheikh Anta Diop de Dakar}
\date{}

\begin{document}

\maketitle

\begin{infobox}{Environnement de travail}
Ce TP doit être réalisé dans l'environnement Docker du cours. 
Si vous n'avez pas encore configuré l'environnement, consultez le TP-1 (Setup Environnement).

\textbf{Commandes de base:}
\begin{verbatim}
# Linux/macOS
docker run -it -v $(pwd):/workspace elbachirtoure/c-advanced-course:v2.0

# Windows (Git Bash)
docker run -it -v /${PWD}:/workspace elbachirtoure/c-advanced-course:v2.0

# Dans le container, aller dans le workspace
cd /workspace
\end{verbatim}
\end{infobox}

\section{Exercice 1 : Manipulation de Structures}

\begin{exercicebox}{1.1 - Création et initialisation}
Créer un fichier \texttt{ex1\_struct.c} avec les structures suivantes:

\begin{lstlisting}
typedef struct {
    char name[32];
    int age;
    float height;
} Person;

typedef struct {
    int day, month, year;
} Date;

typedef struct {
    Person person;
    Date birth_date;
    char *email;  // allocation dynamique
} Employee;
\end{lstlisting}

\textbf{Implémenter:}
\begin{enumerate}
\item \texttt{void print\_person(const Person *p)} -- affiche les infos
\item \texttt{Person create\_person(const char *name, int age, float height)}
\item \texttt{Employee *create\_employee(const char *name, int age, float height, int d, int m, int y, const char *email)}
\item \texttt{void free\_employee(Employee *e)}
\end{enumerate}

\textbf{Tester avec:}
\begin{lstlisting}
Employee *e = create_employee("Alice", 30, 1.65, 15, 3, 1994, "alice@test.com");
// afficher, puis libérer
\end{lstlisting}
\end{exercicebox}

\begin{exercicebox}{1.2 - Comparaison de structures}
\begin{lstlisting}
typedef struct {
    int x, y;
} Point;
\end{lstlisting}

\textbf{Implémenter:}
\begin{enumerate}
\item \texttt{int point\_equal(const Point *a, const Point *b)} -- comparaison membre à membre
\item \texttt{int point\_equal\_memcmp(const Point *a, const Point *b)} -- utiliser \texttt{memcmp}
\end{enumerate}

\textbf{Question:} Pourquoi \texttt{memcmp} peut-il donner des résultats incorrects sur certaines structures ?
\end{exercicebox}

\section{Exercice 2 : Analyse Mémoire avec GDB}

\begin{exercicebox}{2.1 - Inspection de structures}
Créer \texttt{ex2\_gdb.c}:

\begin{lstlisting}
#include <stdio.h>
#include <stddef.h>

struct Data {
    char tag;
    int value;
    short code;
    char flags;
};

int main(void) {
    struct Data d = {'A', 0x12345678, 0xABCD, 'F'};
    
    printf("sizeof(struct Data) = %zu\n", sizeof(struct Data));
    printf("offsetof tag   = %zu\n", offsetof(struct Data, tag));
    printf("offsetof value = %zu\n", offsetof(struct Data, value));
    printf("offsetof code  = %zu\n", offsetof(struct Data, code));
    printf("offsetof flags = %zu\n", offsetof(struct Data, flags));
    
    return 0;
}
\end{lstlisting}

\textbf{Compiler et débugger:}
\begin{verbatim}
gcc -g -o ex2 ex2_gdb.c
gdb ./ex2
\end{verbatim}

\textbf{Commandes GDB à exécuter:}
\begin{verbatim}
(gdb) break main
(gdb) run
(gdb) next
(gdb) ptype struct Data
(gdb) print d
(gdb) print sizeof(d)
(gdb) x/16xb &d
(gdb) print &d.tag
(gdb) print &d.value
\end{verbatim}

\textbf{À rendre:}
\begin{itemize}
\item Capture du dump hexadécimal (\texttt{x/16xb \&d})
\item Identification des octets de padding
\item Vérification de l'endianness
\end{itemize}
\end{exercicebox}

\begin{exercicebox}{2.2 - Structures imbriquées}
\begin{lstlisting}
struct Inner {
    char a;
    int b;
};

struct Outer {
    char x;
    struct Inner inner;
    char y;
};

int main(void) {
    struct Outer o = {'X', {'A', 42}, 'Y'};
    return 0;
}
\end{lstlisting}

\textbf{Avec GDB:}
\begin{enumerate}
\item Afficher \texttt{sizeof(struct Outer)}
\item Examiner le layout avec \texttt{x/Nxb \&o}
\item Identifier tous les octets de padding
\item Dessiner le schéma mémoire
\end{enumerate}
\end{exercicebox}

\section{Exercice 3 : Unions et Type Punning}

\begin{exercicebox}{3.1 - Conversion float/int}
\begin{lstlisting}
#include <stdio.h>

union FloatInt {
    float f;
    unsigned int i;
    unsigned char bytes[4];
};

int main(void) {
    union FloatInt u;
    
    // Tester avec différentes valeurs
    float values[] = {1.0f, -1.0f, 0.0f, 3.14159f, 1e10f};
    
    for (int j = 0; j < 5; j++) {
        u.f = values[j];
        printf("f = %g\n", u.f);
        printf("  hex: 0x%08X\n", u.i);
        printf("  bytes: ");
        for (int i = 0; i < 4; i++) {
            printf("%02X ", u.bytes[i]);
        }
        printf("\n\n");
    }
    
    return 0;
}
\end{lstlisting}

\textbf{Analyser:}
\begin{enumerate}
\item L'ordre des octets (little-endian)
\item La structure IEEE 754 (signe, exposant, mantisse)
\item Pourquoi \texttt{1.0f} donne \texttt{0x3F800000}
\end{enumerate}
\end{exercicebox}

\begin{exercicebox}{3.2 - Tagged Union (Variant)}
Implémenter un type variant pouvant contenir \texttt{int}, \texttt{double}, ou \texttt{char*}:

\begin{lstlisting}
typedef enum { VAR_INT, VAR_DOUBLE, VAR_STRING } VarType;

typedef struct {
    VarType type;
    union {
        int i;
        double d;
        char *s;
    } value;
} Variant;

// À implémenter:
Variant var_int(int i);
Variant var_double(double d);
Variant var_string(const char *s);  // alloue une copie
void var_print(const Variant *v);
void var_free(Variant *v);  // libère si string
\end{lstlisting}

\textbf{Tester avec un tableau de Variants de types différents.}
\end{exercicebox}

\section{Exercice 4 : Optimisation de Structures}

\begin{exercicebox}{4.1 - Mesure et optimisation}
\begin{lstlisting}
// Version originale
struct Original {
    char a;
    double b;
    char c;
    int d;
    short e;
    char f;
};

// Votre version optimisée
struct Optimized {
    // réorganiser les membres
};
\end{lstlisting}

\textbf{Implémenter un programme qui:}
\begin{enumerate}
\item Affiche \texttt{sizeof} et les offsets des deux versions
\item Calcule l'économie en octets et pourcentage
\item Vérifie que les données sont identiques après copie
\end{enumerate}
\end{exercicebox}

\begin{exercicebox}{4.2 - Impact sur les tableaux}
\begin{lstlisting}
#define N 1000000

struct Big { /* version non optimisée */ };
struct Small { /* version optimisée */ };

int main(void) {
    printf("Tableau Big:   %zu Mo\n", N * sizeof(struct Big) / (1024*1024));
    printf("Tableau Small: %zu Mo\n", N * sizeof(struct Small) / (1024*1024));
    return 0;
}
\end{lstlisting}

\textbf{Mesurer l'impact mémoire pour 1 million d'éléments.}
\end{exercicebox}

\section{Exercice 5 : Structures Dynamiques}

\begin{exercicebox}{5.1 - Liste chaînée simple}
\begin{lstlisting}
typedef struct Node {
    int data;
    struct Node *next;
} Node;

typedef struct {
    Node *head;
    size_t size;
} List;

// À implémenter:
List *list_create(void);
void list_push_front(List *list, int data);
void list_push_back(List *list, int data);
int list_pop_front(List *list);
void list_print(const List *list);
void list_destroy(List *list);
\end{lstlisting}

\textbf{Vérifier avec Valgrind:}
\begin{verbatim}
valgrind --leak-check=full ./ex5
\end{verbatim}
\end{exercicebox}

\begin{exercicebox}{5.2 - Flexible Array Member}
Implémenter une chaîne de caractères avec sa longueur:

\begin{lstlisting}
typedef struct {
    size_t length;
    size_t capacity;
    char data[];  // flexible array member
} String;

String *string_create(const char *init);
String *string_concat(const String *a, const String *b);
void string_print(const String *s);
void string_destroy(String *s);
\end{lstlisting}
\end{exercicebox}

\section{Projet : Simple Database}

\begin{techbox}{Description}
Implémenter un système de gestion d'enregistrements en mémoire avec les fonctionnalités suivantes:
\begin{itemize}
\item Stockage dynamique de records (tableau redimensionnable)
\item Opérations CRUD (Create, Read, Update, Delete)
\item Recherche par ID ou par nom
\item Affichage formaté
\end{itemize}
\end{techbox}

\begin{exercicebox}{Structure de données}
\begin{lstlisting}
#define MAX_NAME 64
#define MAX_EMAIL 128

typedef struct {
    int id;
    char name[MAX_NAME];
    char email[MAX_EMAIL];
    int age;
    double salary;
} Record;

typedef struct {
    Record *records;
    size_t count;
    size_t capacity;
    int next_id;
} Database;
\end{lstlisting}
\end{exercicebox}

\begin{exercicebox}{API à implémenter}
\begin{lstlisting}
// Création/Destruction
Database *db_create(size_t initial_capacity);
void db_destroy(Database *db);

// CRUD
int db_insert(Database *db, const char *name, const char *email, 
              int age, double salary);  // retourne l'ID
Record *db_find_by_id(Database *db, int id);
Record *db_find_by_name(Database *db, const char *name);
int db_update(Database *db, int id, const char *name, 
              const char *email, int age, double salary);
int db_delete(Database *db, int id);

// Affichage
void db_print_all(const Database *db);
void db_print_record(const Record *r);

// Statistiques
size_t db_count(const Database *db);
double db_average_salary(const Database *db);
double db_average_age(const Database *db);
\end{lstlisting}
\end{exercicebox}

\begin{exercicebox}{Programme de test}
\begin{lstlisting}
int main(void) {
    Database *db = db_create(10);
    
    // Insertions
    db_insert(db, "Alice", "alice@test.com", 30, 45000.0);
    db_insert(db, "Bob", "bob@test.com", 25, 38000.0);
    db_insert(db, "Charlie", "charlie@test.com", 35, 52000.0);
    
    printf("=== Tous les enregistrements ===\n");
    db_print_all(db);
    
    printf("\n=== Recherche par nom: Bob ===\n");
    Record *r = db_find_by_name(db, "Bob");
    if (r) db_print_record(r);
    
    printf("\n=== Mise a jour ID 2 ===\n");
    db_update(db, 2, "Robert", "robert@test.com", 26, 40000.0);
    db_print_all(db);
    
    printf("\n=== Statistiques ===\n");
    printf("Nombre: %zu\n", db_count(db));
    printf("Salaire moyen: %.2f\n", db_average_salary(db));
    printf("Age moyen: %.1f\n", db_average_age(db));
    
    printf("\n=== Suppression ID 1 ===\n");
    db_delete(db, 1);
    db_print_all(db);
    
    db_destroy(db);
    return 0;
}
\end{lstlisting}
\end{exercicebox}

\begin{warningbox}{Contraintes}
\begin{itemize}
\item Le tableau \texttt{records} doit se redimensionner automatiquement (doubler la capacité)
\item Aucune fuite mémoire (vérifier avec Valgrind)
\item \texttt{db\_delete} peut réorganiser le tableau (compacter)
\item Les ID sont uniques et ne sont jamais réutilisés
\item Gestion des erreurs (NULL, ID invalide, etc.)
\end{itemize}
\end{warningbox}

\section*{Soumission}

\begin{infobox}{Fichiers à rendre}
\begin{itemize}
\item \texttt{ex1\_struct.c} -- Exercice 1
\item \texttt{ex2\_gdb.c} + rapport GDB (captures/schémas)
\item \texttt{ex3\_union.c} -- Exercice 3
\item \texttt{ex4\_optim.c} -- Exercice 4
\item \texttt{ex5\_list.c} -- Exercice 5
\item \texttt{database.h} + \texttt{database.c} + \texttt{main.c} -- Projet
\item \texttt{Makefile}
\item Rapport Valgrind (capture sans erreurs)
\end{itemize}

\textbf{Compilation requise:}
\begin{verbatim}
gcc -Wall -Wextra -Werror -std=c11 -g -o prog *.c
\end{verbatim}
\end{infobox}

\end{document}
